% Generated draft: requires mathematical and source review.
\documentclass[11pt]{article}
\usepackage[margin=25mm]{geometry}
\usepackage{fontspec}
\setmainfont{Latin Modern Roman}
\usepackage{amsmath,amssymb,mathtools,unicode-math}
\setmathfont{Latin Modern Math}
\usepackage{enumitem,needspace,xurl,hyperref}
\setlength{\parindent}{0pt}
\setlength{\parskip}{5pt}
\setlength{\emergencystretch}{3em}
\newcommand{\N}{\mathbb{N}}
\newcommand{\Z}{\mathbb{Z}}
\newcommand{\Q}{\mathbb{Q}}
\newcommand{\R}{\mathbb{R}}
\newcommand{\C}{\mathbb{C}}
\newcommand{\F}{\mathbb{F}}
\newcommand{\A}{\mathbb{A}}
\newcommand{\PP}{\mathbb{P}}
\newcommand{\E}{\mathbb{E}}
\newcommand{\eps}{\varepsilon}
\newcommand{\dd}{\,\mathrm{d}}
\begin{document}
\section*{UnsolvedMath ID: 30002856}
\textbf{Source problem code:} catalogue entry
\subsection{Definitions and mathematical statement}
For a finite group \(K\), let \(\operatorname{Irr}(K)\) denote the set of irreducible complex characters of \(K\), and define its set of character degrees by
\[
\operatorname{cd}(K)=\{\chi(1):\chi\in\operatorname{Irr}(K)\}.
\]
Here \(\chi(1)\) is the dimension of the corresponding irreducible complex representation; degrees are recorded without multiplicity. A nonabelian simple group is a nonabelian group whose only normal subgroups are the identity subgroup and the whole group.

The finite-group formulation suggested by the supplied workshop context and literature assessment is the following conjecture, identified in the supplied assessment as Huppert's character-degree conjecture: for every pair of finite groups \(G,H\), if \(H\) is nonabelian simple and
\[
\operatorname{cd}(G)=\operatorname{cd}(H),
\]
then there exists a finite abelian group \(A\) such that
\[
G\cong H\times A.
\]
The trivial group is allowed for \(A\).

The displayed source statement does not explicitly require \(G\) or \(H\) to be finite and does not define \(\operatorname{cd}\). Thus finiteness and the interpretation as a set of complex character degrees are contextual interpretations requiring confirmation against the cited report; no formulation for infinite groups is supplied.
\subsection{Short English statement}
The intended finite-group conjecture says that a finite group having exactly the same irreducible complex character degrees as a nonabelian finite simple group must be the direct product of that simple group with a finite abelian group. The finite domains and the meaning of character degrees require source confirmation.
\subsection{Sources}
\begin{itemize}
\item Supplied local UnsolvedMath registry record, permanent ID \(30002856\), title ``Character-Degree Rigidity of Nonabelian Simple Groups,'' exact source problem code \texttt{catalogue entry}, in the collection ``Oberwolfach Reports Open Problems.'' The proposer field is empty; the supplied literature assessment identifies the conjecture with Huppert.
\item Supplied source citation: \emph{Representations of Finite Groups (2016)}, Oberwolfach Reports, DOI \texttt{10.4171/owr/2015/18}, pp.\ 971--1027. Source URL and literature-source label ``Oberwolfach Report 18'': \url{https://doi.org/10.4171/owr/2015/18}. The record gives proposed year 2015; its report title contains 2016. These supplied bibliographic details have not been independently reconciled.
\item Supplied literature source: Alavi--Daneshkhah--Tong-Viet--Wakefield, \url{https://doi.org/10.1017/S1446788712000535}.
\item Supplied literature source: Tong-Viet, \emph{Journal of Algebra} 2025, \url{https://doi.org/10.1016/j.jalgebra.2024.10.042}.
\item Literature assessment dated exactly \texttt{2026-08-21}: the supplied curation labels the problem ``partially solved'' and reports proofs for several major families, including alternating, sporadic, and exceptional Lie-type groups, while describing the general finite-simple-group statement as unresolved. This is supplied dated curation, not an independently verified literature assessment. The source background explicitly requests independent verification of time-sensitive status and cited claims.
\item The registry labels the statement \texttt{corrected\_verified} and states that mathematical content was preserved while notation, typography, or omitted definitions were repaired for publication. This is a supplied editorial qualification, not independent mathematical verification of this draft.
\end{itemize}
\end{document}
